% GENERATED FILE. Edit canonical lesson frontmatter, then run npm run generate:books.
% canonical-entries: 92

\chapter*{Glossary}
\addcontentsline{toc}{chapter}{Glossary}
\markboth{Glossary}{Glossary}

This glossary collects every word and phrase taught in the book. Pronunciation appears when it differs from the written form; chapter numbers show where each entry is introduced.

\noindent\begin{minipage}[t]{\linewidth}
\raggedright
\textbf{a água, o vinho}\par
\small water and wine — água keeps aqua's shape; vinho shows the -nh- for Latin -n-\par
\footnotesize Introduced in Chapter 11.
\end{minipage}\par\medskip

\noindent\begin{minipage}[t]{\linewidth}
\raggedright
\textbf{a boca}\par
\small mouth — Vulgar Latin's word for \textquotedblleft{}cheek,\textquotedblright{} promoted to cover the whole mouth\par
\footnotesize Introduced in Chapter 26.
\end{minipage}\par\medskip

\noindent\begin{minipage}[t]{\linewidth}
\raggedright
\textbf{a cabeça}\par
\small the head — feminine, with cedilla ç pronounced s\par
\footnotesize Introduced in Chapter 17.
\end{minipage}\par\medskip

\noindent\begin{minipage}[t]{\linewidth}
\raggedright
\textbf{a família}\par
\small family — a word that used to mean the household's servants, not its blood relatives\par
\footnotesize Introduced in Chapter 25.
\end{minipage}\par\medskip

\noindent\begin{minipage}[t]{\linewidth}
\raggedright
\textbf{a mão}\par
\small the hand — feminine, and a showcase of the -ão ending that swallowed three Latin endings\par
\footnotesize Introduced in Chapter 17.
\end{minipage}\par\medskip

\noindent\begin{minipage}[t]{\linewidth}
\raggedright
\textbf{a mulher}\par
\small woman — the fourth Romance sister to pick a fourth different Latin word\par
\footnotesize Introduced in Chapter 25.
\end{minipage}\par\medskip

\noindent\begin{minipage}[t]{\linewidth}
\raggedright
\textbf{a orelha}\par
\small ear — a second word built by the same -cl- to -lh- law, and the same PT/ES split as trabalhar\par
\footnotesize Introduced in Chapter 26.
\end{minipage}\par\medskip

\noindent\begin{minipage}[t]{\linewidth}
\raggedright
\textbf{adeus}\par
\small goodbye (literally \textquotedblleft{}to God\textquotedblright{})\par
\footnotesize Introduced in Chapter 4.
\end{minipage}\par\medskip

\noindent\begin{minipage}[t]{\linewidth}
\raggedright
\textbf{ajudar}\par
\small to help — and the Latin i that hardened into Portuguese j\par
\footnotesize Introduced in Chapter 20.
\end{minipage}\par\medskip

\noindent\begin{minipage}[t]{\linewidth}
\raggedright
\textbf{as estações}\par
\small the seasons — and verão, a \textquotedblleft{}summer\textquotedblright{} that came from \textquotedblleft{}spring\textquotedblright{}\par
\footnotesize Introduced in Chapter 9.
\end{minipage}\par\medskip

\noindent\begin{minipage}[t]{\linewidth}
\raggedright
\textbf{até amanhã}\par
\small see you tomorrow (literally \textquotedblleft{}until tomorrow\textquotedblright{})\par
\footnotesize Introduced in Chapter 4.
\end{minipage}\par\medskip

\noindent\begin{minipage}[t]{\linewidth}
\raggedright
\textbf{até breve}\par
\small see you soon (literally \textquotedblleft{}until {[}a{]} short {[}while{]}\textquotedblright{})\par
\footnotesize Introduced in Chapter 4.
\end{minipage}\par\medskip

\noindent\begin{minipage}[t]{\linewidth}
\raggedright
\textbf{até logo}\par
\small see you later (literally \textquotedblleft{}until soon/then\textquotedblright{})\par
\footnotesize Introduced in Chapter 4.
\end{minipage}\par\medskip

\noindent\begin{minipage}[t]{\linewidth}
\raggedright
\textbf{bom / boa}\par
\small good (adjective)\par
\footnotesize Introduced in Chapter 1.
\end{minipage}\par\medskip

\noindent\begin{minipage}[t]{\linewidth}
\raggedright
\textbf{bom dia}\par
\small good morning / good day\par
\footnotesize Introduced in Chapter 1.
\end{minipage}\par\medskip

\noindent\begin{minipage}[t]{\linewidth}
\raggedright
\textbf{como}\par
\small how / like, as\par
\footnotesize Introduced in Chapter 2.
\end{minipage}\par\medskip

\noindent\begin{minipage}[t]{\linewidth}
\raggedright
\textbf{Como se chama?}\par
\small what's your name? (literally \textquotedblleft{}how do you call yourself?\textquotedblright{})\par
\footnotesize Introduced in Chapter 3.
\end{minipage}\par\medskip

\noindent\begin{minipage}[t]{\linewidth}
\raggedright
\textbf{Como vai? / Como está?}\par
\small how are you? — the “go” and “stand” versions\par
\footnotesize Introduced in Chapter 2.
\end{minipage}\par\medskip

\noindent\begin{minipage}[t]{\linewidth}
\raggedright
\textbf{comprar}\par
\small to buy — from a Latin verb spelled exactly like the one behind 'compare', and not the same verb at all\par
\footnotesize Introduced in Chapter 21.
\end{minipage}\par\medskip

\noindent\begin{minipage}[t]{\linewidth}
\raggedright
\textbf{conseguir / obter}\par
\small to get — the everyday verb that also means 'to manage to', and the formal twin built on ter\par
\footnotesize Introduced in Chapter 21.
\end{minipage}\par\medskip

\noindent\begin{minipage}[t]{\linewidth}
\raggedright
\textbf{de nada}\par
\small you're welcome (literally \textquotedblleft{}of nothing\textquotedblright{})\par
\footnotesize Introduced in Chapter 2.
\end{minipage}\par\medskip

\noindent\begin{minipage}[t]{\linewidth}
\raggedright
\textbf{dezesseis — vinte}\par
\small 16–20 — Portuguese breaks earliest, and rebuilds numbers as \textquotedblleft{}ten AND six\textquotedblright{}\par
\footnotesize Introduced in Chapter 12.
\end{minipage}\par\medskip

\noindent\begin{minipage}[t]{\linewidth}
\raggedright
\textbf{dia}\par
\small day (masculine)\par
\footnotesize Introduced in Chapter 1.
\end{minipage}\par\medskip

\noindent\begin{minipage}[t]{\linewidth}
\raggedright
\textbf{dizer}\par
\small to say — and the Latin verb that also produced the judge\par
\footnotesize Introduced in Chapter 18.
\end{minipage}\par\medskip

\noindent\begin{minipage}[t]{\linewidth}
\raggedright
\textbf{encontrar}\par
\small to find\par
\footnotesize Introduced in Chapter 22.
\end{minipage}\par\medskip

\noindent\begin{minipage}[t]{\linewidth}
\raggedright
\textbf{encontrar-se}\par
\small to meet up\par
\footnotesize Introduced in Chapter 22.
\end{minipage}\par\medskip

\noindent\begin{minipage}[t]{\linewidth}
\raggedright
\textbf{entender / compreender}\par
\small to understand — one verb that stretches toward the thing, one that seizes it whole\par
\footnotesize Introduced in Chapter 19.
\end{minipage}\par\medskip

\noindent\begin{minipage}[t]{\linewidth}
\raggedright
\textbf{escrever}\par
\small to write — to scratch, and the propped-up e- that decodes a hundred other words\par
\footnotesize Introduced in Chapter 19.
\end{minipage}\par\medskip

\noindent\begin{minipage}[t]{\linewidth}
\raggedright
\textbf{esperar}\par
\small to expect\par
\footnotesize Introduced in Chapter 21.
\end{minipage}\par\medskip

\noindent\begin{minipage}[t]{\linewidth}
\raggedright
\textbf{esperar}\par
\small to wait\par
\footnotesize Introduced in Chapter 21.
\end{minipage}\par\medskip

\noindent\begin{minipage}[t]{\linewidth}
\raggedright
\textbf{espero que sim}\par
\small I hope so\par
\footnotesize Introduced in Chapter 21.
\end{minipage}\par\medskip

\noindent\begin{minipage}[t]{\linewidth}
\raggedright
\textbf{eu}\par
\small I (first-person subject pronoun)\par
\footnotesize Introduced in Chapter 3.
\end{minipage}\par\medskip

\noindent\begin{minipage}[t]{\linewidth}
\raggedright
\textbf{falar}\par
\small to speak / to talk (your first regular -ar verb)\par
\footnotesize Introduced in Chapter 5.
\end{minipage}\par\medskip

\noindent\begin{minipage}[t]{\linewidth}
\raggedright
\textbf{falei}\par
\small the pretérito perfeito — Portuguese kept its simple past, where three sisters gave theirs up\par
\footnotesize Introduced in Chapter 15.
\end{minipage}\par\medskip

\noindent\begin{minipage}[t]{\linewidth}
\raggedright
\textbf{Falo português}\par
\small \textquotedblleft{}I speak Portuguese\textquotedblright{} — your first full sentence (português = Portuguese)\par
\footnotesize Introduced in Chapter 5.
\end{minipage}\par\medskip

\noindent\begin{minipage}[t]{\linewidth}
\raggedright
\textbf{gostar de}\par
\small to like, to love — the verb that will not let go of its preposition\par
\footnotesize Introduced in Chapter 20.
\end{minipage}\par\medskip

\noindent\begin{minipage}[t]{\linewidth}
\raggedright
\textbf{há três anos}\par
\small three years ago\par
\footnotesize Introduced in Chapter 18.
\end{minipage}\par\medskip

\noindent\begin{minipage}[t]{\linewidth}
\raggedright
\textbf{hora}\par
\small hour / o'clock — telling the time\par
\footnotesize Introduced in Chapter 8.
\end{minipage}\par\medskip

\noindent\begin{minipage}[t]{\linewidth}
\raggedright
\textbf{ir}\par
\small to go — two letters long, and welded together out of three separate Latin verbs\par
\footnotesize Introduced in Chapter 18.
\end{minipage}\par\medskip

\noindent\begin{minipage}[t]{\linewidth}
\raggedright
\textbf{jogar / brincar / tocar}\par
\small to play — split three ways, by whether it is a game, a child at play, or an instrument\par
\footnotesize Introduced in Chapter 22.
\end{minipage}\par\medskip

\noindent\begin{minipage}[t]{\linewidth}
\raggedright
\textbf{ler}\par
\small to read — three letters, irregular, and the same broken row as ver\par
\footnotesize Introduced in Chapter 19.
\end{minipage}\par\medskip

\noindent\begin{minipage}[t]{\linewidth}
\raggedright
\textbf{mais ou menos}\par
\small so-so (literally \textquotedblleft{}more or less\textquotedblright{})\par
\footnotesize Introduced in Chapter 2.
\end{minipage}\par\medskip

\noindent\begin{minipage}[t]{\linewidth}
\raggedright
\textbf{me chamo / chamo-me}\par
\small my name is (literally \textquotedblleft{}I call myself\textquotedblright{})\par
\footnotesize Introduced in Chapter 3.
\end{minipage}\par\medskip

\noindent\begin{minipage}[t]{\linewidth}
\raggedright
\textbf{meio-dia, meia-noite}\par
\small noon and midnight — \textquotedblleft{}mid-day\textquotedblright{} and \textquotedblleft{}mid-night\textquotedblright{}\par
\footnotesize Introduced in Chapter 8.
\end{minipage}\par\medskip

\noindent\begin{minipage}[t]{\linewidth}
\raggedright
\textbf{morar}\par
\small to live / to dwell (a second -ar verb — from \textquotedblleft{}to linger\textquotedblright{})\par
\footnotesize Introduced in Chapter 5.
\end{minipage}\par\medskip

\noindent\begin{minipage}[t]{\linewidth}
\raggedright
\textbf{noite / boa noite}\par
\small night / good night\par
\footnotesize Introduced in Chapter 1.
\end{minipage}\par\medskip

\noindent\begin{minipage}[t]{\linewidth}
\raggedright
\textbf{o / a}\par
\small the (masculine / feminine)\par
\footnotesize Introduced in Chapter 1.
\end{minipage}\par\medskip

\noindent\begin{minipage}[t]{\linewidth}
\raggedright
\textbf{o açúcar}\par
\small sugar — the same route as arroz, and the same cedilla as cabeça\par
\footnotesize Introduced in Chapter 24.
\end{minipage}\par\medskip

\noindent\begin{minipage}[t]{\linewidth}
\raggedright
\textbf{o amigo, a amiga}\par
\small friend — from the verb \textquotedblleft{}to love,\textquotedblright{} and the word you close an introduction with\par
\footnotesize Introduced in Chapter 25.
\end{minipage}\par\medskip

\noindent\begin{minipage}[t]{\linewidth}
\raggedright
\textbf{o arroz}\par
\small rice — Greek by way of Arabic, the al-Andalus thread again\par
\footnotesize Introduced in Chapter 24.
\end{minipage}\par\medskip

\noindent\begin{minipage}[t]{\linewidth}
\raggedright
\textbf{o café}\par
\small coffee — a word Portuguese never inherited from Latin at all\par
\footnotesize Introduced in Chapter 23.
\end{minipage}\par\medskip

\noindent\begin{minipage}[t]{\linewidth}
\raggedright
\textbf{o chá}\par
\small tea — Portuguese is the ONE Western European exception in how it got this word\par
\footnotesize Introduced in Chapter 23.
\end{minipage}\par\medskip

\noindent\begin{minipage}[t]{\linewidth}
\raggedright
\textbf{o coração}\par
\small heart — the same nasal -ão ending as pão, but GROWN rather than worn down\par
\footnotesize Introduced in Chapter 26.
\end{minipage}\par\medskip

\noindent\begin{minipage}[t]{\linewidth}
\raggedright
\textbf{o homem}\par
\small man — from the Latin word for \textquotedblleft{}human being,\textquotedblright{} narrowed to just one sex\par
\footnotesize Introduced in Chapter 25.
\end{minipage}\par\medskip

\noindent\begin{minipage}[t]{\linewidth}
\raggedright
\textbf{o irmão, a irmã}\par
\small brother and sister — NOT from frater, but from germanus \textquotedblleft{}of the same stock\textquotedblright{}\par
\footnotesize Introduced in Chapter 10.
\end{minipage}\par\medskip

\noindent\begin{minipage}[t]{\linewidth}
\raggedright
\textbf{o leite}\par
\small milk — the one drink of the three that IS plain inherited Latin\par
\footnotesize Introduced in Chapter 23.
\end{minipage}\par\medskip

\noindent\begin{minipage}[t]{\linewidth}
\raggedright
\textbf{o olho}\par
\small eye — the same -cl- to -lh- sound law that gave Chapter 13 its \textquotedblleft{}worm\textquotedblright{} red\par
\footnotesize Introduced in Chapter 26.
\end{minipage}\par\medskip

\noindent\begin{minipage}[t]{\linewidth}
\raggedright
\textbf{o ovo}\par
\small egg — three letters, unchanged for two thousand years\par
\footnotesize Introduced in Chapter 24.
\end{minipage}\par\medskip

\noindent\begin{minipage}[t]{\linewidth}
\raggedright
\textbf{o pai, a mãe}\par
\small father and mother — the most worn-down of all: pater$\to$pai, mater$\to$mãe\par
\footnotesize Introduced in Chapter 10.
\end{minipage}\par\medskip

\noindent\begin{minipage}[t]{\linewidth}
\raggedright
\textbf{o pão}\par
\small bread — pānis worn to a single nasal syllable, just like pai\par
\footnotesize Introduced in Chapter 11.
\end{minipage}\par\medskip

\noindent\begin{minipage}[t]{\linewidth}
\raggedright
\textbf{o queijo}\par
\small cheese — the word English and Portuguese still share, though French and Italian moved on\par
\footnotesize Introduced in Chapter 24.
\end{minipage}\par\medskip

\noindent\begin{minipage}[t]{\linewidth}
\raggedright
\textbf{obrigado / obrigada}\par
\small thank you\par
\footnotesize Introduced in Chapter 1.
\end{minipage}\par\medskip

\noindent\begin{minipage}[t]{\linewidth}
\raggedright
\textbf{olá}\par
\small hi\par
\footnotesize Introduced in Chapter 1.
\end{minipage}\par\medskip

\noindent\begin{minipage}[t]{\linewidth}
\raggedright
\textbf{onze — quinze}\par
\small 11–15 — the only five teens Portuguese kept fused; it breaks earliest of all\par
\footnotesize Introduced in Chapter 12.
\end{minipage}\par\medskip

\noindent\begin{minipage}[t]{\linewidth}
\raggedright
\textbf{os meses}\par
\small the months — Roman gods, emperors, and the numbers 7–10\par
\footnotesize Introduced in Chapter 9.
\end{minipage}\par\medskip

\noindent\begin{minipage}[t]{\linewidth}
\raggedright
\textbf{pensar}\par
\small to think — from a Latin verb that meant to weigh a thing in a balance\par
\footnotesize Introduced in Chapter 19.
\end{minipage}\par\medskip

\noindent\begin{minipage}[t]{\linewidth}
\raggedright
\textbf{perguntar}\par
\small to ask a question — and why it is not the verb for asking for a coffee\par
\footnotesize Introduced in Chapter 20.
\end{minipage}\par\medskip

\noindent\begin{minipage}[t]{\linewidth}
\raggedright
\textbf{prazer}\par
\small pleased to meet you (literally \textquotedblleft{}a pleasure\textquotedblright{})\par
\footnotesize Introduced in Chapter 3.
\end{minipage}\par\medskip

\noindent\begin{minipage}[t]{\linewidth}
\raggedright
\textbf{preto, branco}\par
\small black and white — Portuguese has TWO blacks, and its white is a Germanic loan\par
\footnotesize Introduced in Chapter 13.
\end{minipage}\par\medskip

\noindent\begin{minipage}[t]{\linewidth}
\raggedright
\textbf{responder}\par
\small to answer — literally to pledge something back, and the other half of asking\par
\footnotesize Introduced in Chapter 22.
\end{minipage}\par\medskip

\noindent\begin{minipage}[t]{\linewidth}
\raggedright
\textbf{sábado, domingo}\par
\small the weekend (Saturday, Sunday) — the two days that kept their names\par
\footnotesize Introduced in Chapter 7.
\end{minipage}\par\medskip

\noindent\begin{minipage}[t]{\linewidth}
\raggedright
\textbf{saber / conhecer}\par
\small to know — the second place Portuguese answers one English verb with two\par
\footnotesize Introduced in Chapter 18.
\end{minipage}\par\medskip

\noindent\begin{minipage}[t]{\linewidth}
\raggedright
\textbf{segunda, terça, quarta, quinta, sexta(-feira)}\par
\small the weekdays (Mon–Fri) — Portuguese NUMBERS its days instead of naming planets\par
\footnotesize Introduced in Chapter 7.
\end{minipage}\par\medskip

\noindent\begin{minipage}[t]{\linewidth}
\raggedright
\textbf{seis, sete, oito, nove, dez}\par
\small the numbers 6–10 (and the months hidden in them)\par
\footnotesize Introduced in Chapter 6.
\end{minipage}\par\medskip

\noindent\begin{minipage}[t]{\linewidth}
\raggedright
\textbf{ser}\par
\small to be — present, imperfect, and preterite forms\par
\footnotesize Introduced in Chapter 16.
\end{minipage}\par\medskip

\noindent\begin{minipage}[t]{\linewidth}
\raggedright
\textbf{ser / estar}\par
\small to be — the one English verb that Portuguese answers with two, and the line it draws its own way\par
\footnotesize Introduced in Chapter 18.
\end{minipage}\par\medskip

\noindent\begin{minipage}[t]{\linewidth}
\raggedright
\textbf{ser vs. estar}\par
\small identity or nature against location or resulting condition\par
\footnotesize Introduced in Chapter 16.
\end{minipage}\par\medskip

\noindent\begin{minipage}[t]{\linewidth}
\raggedright
\textbf{tarde / boa tarde}\par
\small afternoon / good afternoon\par
\footnotesize Introduced in Chapter 1.
\end{minipage}\par\medskip

\noindent\begin{minipage}[t]{\linewidth}
\raggedright
\textbf{tenho falado}\par
\small ter + participle — which does NOT mean 'I have spoken', the trap English speakers walk into\par
\footnotesize Introduced in Chapter 15.
\end{minipage}\par\medskip

\noindent\begin{minipage}[t]{\linewidth}
\raggedright
\textbf{tenho vinte anos}\par
\small age — Portuguese HOLDS its years, and annus loses its double n\par
\footnotesize Introduced in Chapter 14.
\end{minipage}\par\medskip

\noindent\begin{minipage}[t]{\linewidth}
\raggedright
\textbf{ter}\par
\small to have — where Portuguese and Spanish broke ranks with the rest of Romance\par
\footnotesize Introduced in Chapter 14.
\end{minipage}\par\medskip

\noindent\begin{minipage}[t]{\linewidth}
\raggedright
\textbf{ter / haver}\par
\small to have, and there is\par
\footnotesize Introduced in Chapter 18.
\end{minipage}\par\medskip

\noindent\begin{minipage}[t]{\linewidth}
\raggedright
\textbf{ter que}\par
\small to have to\par
\footnotesize Introduced in Chapter 18.
\end{minipage}\par\medskip

\noindent\begin{minipage}[t]{\linewidth}
\raggedright
\textbf{tomar / pegar}\par
\small to take — split by whether the thing ends up in your hand or inside you\par
\footnotesize Introduced in Chapter 20.
\end{minipage}\par\medskip

\noindent\begin{minipage}[t]{\linewidth}
\raggedright
\textbf{trabalhar}\par
\small to work (a third -ar verb — back on the \textquotedblleft{}torture\textquotedblright{} road)\par
\footnotesize Introduced in Chapter 5.
\end{minipage}\par\medskip

\noindent\begin{minipage}[t]{\linewidth}
\raggedright
\textbf{trazer}\par
\small to bring — carrying a thing toward the person speaking, from a Latin verb that meant to drag\par
\footnotesize Introduced in Chapter 21.
\end{minipage}\par\medskip

\noindent\begin{minipage}[t]{\linewidth}
\raggedright
\textbf{tudo}\par
\small everything / all\par
\footnotesize Introduced in Chapter 2.
\end{minipage}\par\medskip

\noindent\begin{minipage}[t]{\linewidth}
\raggedright
\textbf{Tudo bem?}\par
\small everything well? — the verb-free everyday “how are you?”\par
\footnotesize Introduced in Chapter 2.
\end{minipage}\par\medskip

\noindent\begin{minipage}[t]{\linewidth}
\raggedright
\textbf{um, dois, três, quatro, cinco}\par
\small the numbers 1–5\par
\footnotesize Introduced in Chapter 6.
\end{minipage}\par\medskip

\noindent\begin{minipage}[t]{\linewidth}
\raggedright
\textbf{ver}\par
\small to see — and the one form that belongs to two different verbs at once\par
\footnotesize Introduced in Chapter 18.
\end{minipage}\par\medskip

\noindent\begin{minipage}[t]{\linewidth}
\raggedright
\textbf{vermelho, azul}\par
\small red and blue — a red named after a worm, and a blue named after a stone\par
\footnotesize Introduced in Chapter 13.
\end{minipage}\par\medskip

\noindent\begin{minipage}[t]{\linewidth}
\raggedright
\textbf{vir}\par
\small to come — and the missing n that explains why the infinitive is so short\par
\footnotesize Introduced in Chapter 18.
\end{minipage}\par\medskip
